% BHU PYQ -- Manually transcribed & error-corrected
\documentclass[11pt,a4paper]{article}

\usepackage[a4paper,top=2.5cm,bottom=2.5cm,left=2.8cm,right=2.8cm,headheight=14pt]{geometry}
\usepackage{amsmath,amssymb}
\usepackage[T1]{fontenc}
\usepackage[utf8]{inputenc}
\usepackage{lmodern}
\usepackage{microtype}
\usepackage{fancyhdr}
\usepackage{enumitem}
\usepackage{hyperref}
\usepackage{lastpage}
\usepackage{parskip}

\hypersetup{colorlinks=true,urlcolor=black,linkcolor=black,
  pdftitle={CHB-504: Physical Chemistry-III -- BHU PYQ},pdfauthor={Banaras Hindu University}}

\pagestyle{fancy}\fancyhf{}
\renewcommand{\headrulewidth}{0.4pt}\renewcommand{\footrulewidth}{0.4pt}
\fancyhead[L]{\small   \textit{Banaras Hindu University}}
\fancyhead[R]{\small   \textit{CHB-504: Physical Chemistry-III}}
\fancyfoot[C]{\small Page  \thepage\ of \pageref{LastPage}}


\newlist{parts}{enumerate}{2}
\setlist[parts,1]{label=(\alph*),leftmargin=*,itemsep=5pt,topsep=3pt}
\setlist[parts,2]{label=(\roman*),leftmargin=*,itemsep=3pt,topsep=2pt}

\newcommand{\pts}[1]{\hfill   \textit{[#1]}}


\begin{document}


\begin{center}
  {\large\bfseries BANARAS HINDU UNIVERSITY}\\[2pt]
  {\normalsize B.Sc. (Hons.) Semester V Examination, 2013-14}\\[6pt]
  \rule{0.8\linewidth}{0.5pt}\\[6pt]
  {\large\bfseries Paper No. CHB-504: Physical Chemistry-III}\\[4pt]
  {\normalsize Time: 3 Hours\hfill Maximum Marks: 70}
\end{center}
\rule{\linewidth}{0.4pt}
\small



\noindent   \textit{Instructions: Answer five questions in total, including Question~1 which is compulsory.}

\normalsize
\bigskip




\noindent   \textbf{Question 1.} Answer all parts of the following: \pts{5 $ \times$ 2 = 10}

\begin{parts}
  \item Why does a solvent flow spontaneously into a solution when separated by a semi-permeable membrane?
  \item Differentiate between activity ($a$) and activity coefficient ($\gamma$).
  \item Write the Debye-H\"uckel-Onsager equation for equivalent conductance.
  \item What is the thickness (or radius) of the ionic cloud surrounding a reference ion?
  \item What are the primary radiolytic products of water?
\end{parts}

\medskip\rule{\linewidth}{0.2pt}




\noindent   \textbf{Question 2.}

\begin{parts}
  \item Discuss the concept and physical significance of activity and activity coefficient. Outline the conditions under which reference and standard states are defined for non-ideal solutions. \pts{6}
  \item State Raoult's law and Henry's law. Discuss the thermodynamic properties of ideal-dilute solutions. \pts{8}
\end{parts}

\medskip\rule{\linewidth}{0.2pt}




\noindent   \textbf{Question 3.}

\begin{parts}
  \item Give a qualitative idea of the Debye-H\"uckel theory of ion-ion interactions and show how these interactions are accounted for. \pts{8}
  \item Derive the relation between mean ionic activity coefficient ($\gamma_\pm$) and ionic strength ($I$) for dilute solutions. \pts{6}
\end{parts}

\medskip\rule{\linewidth}{0.2pt}




\noindent   \textbf{Question 4.}

\begin{parts}
  \item Discuss the different modes of interaction of $\gamma$-rays with matter (photoelectric effect, Compton scattering, and pair production). \pts{6}
  \item Discuss the primary radiolytic products and secondary chemical reactions in the radiolysis of an aqueous solution of ferrous sulphate ($ \text{FeSO}_4$). \pts{8}
\end{parts}

\medskip\rule{\linewidth}{0.2pt}




\noindent   \textbf{Question 5.} Write short notes on any two of the following: \pts{2 $ \times$ 7 = 14}

\begin{parts}
  \item Bravais lattices and crystal systems
  \item Michaelis-Menten kinetics for enzyme catalysis
  \item Nuclear fission and fusion processes
\end{parts}

\vfill
\rule{\linewidth}{0.4pt}

\begin{center}\small
  \href{https://bhu-exam.vercel.app/}{Click here to visit our website}\\[4pt]
  \href{https://me-aryan.vercel.app/}{Click here to visit our contributor}\\[4pt]
  \href{https://quashan-me.vercel.app/}{Click here to visit our contributor}
\end{center}

\end{document}
